\documentclass[letterpaper]{article}
\usepackage[submission]{aaai2027}
\usepackage[hyphens]{url}
\usepackage{graphicx}
\urlstyle{rm}
\def\UrlFont{\rm}
\usepackage{natbib}
\usepackage{booktabs}
\usepackage{amsmath,amssymb}
\frenchspacing

\pdfinfo{
/TemplateVersion (2027.1)
}

\setcounter{secnumdepth}{0}

\title{Supplementary Material for ARMS: Routing by Model Capabilities for
Evolving Vision-Language Model Pools}
\author{Anonymous Submission}
\affiliations{}

\begin{document}
\maketitle

\section{Scope}

This supplement provides additional implementation and evaluation details for
ARMS, including routing supervision, capability descriptors, optimization,
pool evolution, and ablation settings.

\section{Routing Supervision Protocol}

Let $D_{\rm old}$ denote the labeled query--candidate pairs used to train the
router for the original four candidates, and $D_{\rm new}$ the pairs for the
three newly added candidates used in the evolved-pool experiments. The two
annotation groups are maintained separately. The ID/OOD sources and
partitioning follow the main paper.

For query $i$, candidate $j$, candidate response $\hat y_{i,j}$, and reference
answer $y_i$, the suitability label is
\begin{equation}
\bar y_{i,j}=\mathbb{1}\!\left[
\operatorname{norm}(\hat y_{i,j})=
\operatorname{norm}(y_i)\right].
\label{eq:label}
\end{equation}
For each candidate pool, we retain query $i$ only if
$\max_j\bar y_{i,j}=1$; that is, at least one candidate is correct under
Equation~\ref{eq:label}. The oracle therefore obtains EM 1.0 by construction.
This protocol evaluates candidate selection conditional on a correct candidate
being available. Queries on which all candidates fail are excluded from both
training and evaluation; failure detection and abstention are not evaluated.
All outcomes for one query remain in the same partition; splitting is performed
at the query level rather than the query--candidate-pair level.

\paragraph{Prompting and parsing.}
Multiple-choice prompts request only the letter of the correct answer. The
parser extracts the first valid standalone option letter and compares it
case-insensitively with the reference. Free-form normalization strips leading
and trailing whitespace, lowercases, and collapses consecutive whitespace
before strict equality. It does not apply semantic or substring matching.

\section{Candidates and Capability Descriptors}

The original pool contains Qwen2.5-VL-7B, Idefics2-8B, SmolVLM2-2B, and
Llama-3.2-11B-Vision. The evolved pool adds GPT-4o, Claude 3.5 Sonnet, and
Gemini 1.5 Flash.

\paragraph{Candidate inference.}
Open-source candidate responses are generated with bfloat16 weights, each
processor's chat template,
\texttt{max\_new\_tokens=16}, and deterministic greedy decoding
(\texttt{do\_sample=False}).

Each candidate is represented by one structured textual descriptor. The
descriptor records the candidate identity and model version, parameter scale,
architecture family, vision backbone, language backbone, and a concise
summary of documented capabilities. The same schema is used for open-source and
commercial candidates. Descriptors are encoded once and cached because they do
not depend on the query. Upon publication, the companion artifact will provide
the seven verbatim descriptor strings together with provider model identifiers,
access dates for commercial APIs, and immutable checkpoint revisions for
open-source models.

\section{ARMS Implementation}

\subsection{Encoders, fusion, and matching}

The query text and each capability descriptor are encoded with DistilBERT. A
CLIP-ViT-B/32 image encoder produces visual tokens. The question representation
has dimension $d_q=768$, the visual tokens have dimension $d_g=512$, and the
descriptor embedding has dimension $d_m=768$.

A query-dependent gate assigns a soft weight to each of four cross-attention
experts:
\begin{equation}
p_i=\operatorname{softmax}\!\left(
W_{\rm gate}[q_i;\bar g_i]+b_{\rm gate}\right).
\end{equation}
Each expert uses the question as query over visual tokens. If
$a_i^{(n)}$ is expert $n$'s attended visual feature and $e_i^{(n)}$ its fused
output, the shared query representation is
\begin{equation}
e_i=\sum_{n=1}^{4}p_i^{(n)}e_i^{(n)}.
\end{equation}
After concatenating $e_i$ with the question representation, ARMS projects the
result into the descriptor space and computes cosine
similarities. The cosine is clipped away from $\{-1,1\}$ before the bounded
inverse-hyperbolic-tangent transform and sigmoid described in the main paper.
Descriptor embeddings are cached, so query encoding and expert fusion are
shared across all candidates.

\subsection{Candidate-order invariance and complexity}

ARMS applies the same scoring function $f_{\Phi}(x_i,M_j)$ to every candidate.
Permuting the candidate list therefore permutes the scores but does not change
the selected model. Adding a candidate also leaves the learned output shape
unchanged: the new descriptor is encoded into the same matching space instead
of adding a new coordinate to a fixed classification head. This structural
property supports pool evolution, although reliable routing to an unseen model
still requires outcome supervision.

Let $C_{\rm enc}$ be the computational cost of the image and text encoders,
$C_{\rm att}$ the cost of one fusion expert, $K$ the candidate count, and $d_m$
the descriptor dimension. With cached descriptors, per-query complexity is
\begin{equation}
O(C_{\rm enc}+4C_{\rm att}+Kd_m).
\end{equation}
Only the final matching term grows with the pool, and no candidate VLM is invoked
before selection. This expression is asymptotic; latency and throughput are not
measured.

\subsection{Optimization and hardware}

\begin{table}[t]
\centering
\small
\begin{tabular}{ll}
\toprule
Component & Reported setting\\
\midrule
Image encoder & CLIP-ViT-B/32\\
Text/descriptor encoder & DistilBERT\\
Fusion experts & 4\\
Visual/text/descriptor dims. & 512 / 768 / 768\\
Batch size & 64\\
Optimizer & AdamW\\
Learning rate & $8\times10^{-5}$\\
Weight decay & $10^{-2}$\\
Router hardware & $4\times$RTX 4090\\
Open-model output collection & $4\times$A100\\
\bottomrule
\end{tabular}
\caption{Core implementation settings reported for ARMS.}
\label{tab:implementation}
\end{table}

The router hardware also includes a 16-core Xeon Gold 6430 CPU. Binary
cross-entropy is applied to the query--candidate labels in
Equation~\ref{eq:label}. At evaluation time, all candidates for a query are
scored and the candidate with the maximum score is selected.

\section{Protocol for an Evolving VLM Pool}

\paragraph{Incremental training.}
Starting from the original router $\Phi$, incremental training continues
optimization on $D_{\rm new}$ and produces $\Phi+\Delta\Phi$. This reuses the
original representation but can be sensitive to imbalance between supervision
for the original and new candidates.

\paragraph{Independent training.}
A second router $\Theta$ is trained only on $D_{\rm new}$. For query $x$, the
original router first returns its top candidate $V_o$ and confidence $c_o$. If
$c_o>\theta$, $V_o$ is accepted. Otherwise, the new-pool router returns
$(V_n,c_n)$, and ARMS selects $V_o$ when $c_o\geq c_n$ and $V_n$ otherwise.
Thus, independent training avoids parameter interference and invokes the
new-pool router only when the old-pool confidence does not pass the cascade
threshold.

The threshold analysis in the main paper reports the full tested range rather
than relying on OOD labels as an optimization signal. The two router maxima are
compared on the confidence scale defined above, and the reported curve is
interpreted without assuming that the two routers are perfectly calibrated.

\section{Evaluation and Additional Results}

\noindent\textbf{Answer-level and pairwise metrics.}\par
\noindent
Let $j_i^\star=\arg\max_j c_{i,j}$. The primary metric is
\begin{equation}
\operatorname{EM}=\frac{1}{|\mathcal{Q}|}
\sum_{i\in\mathcal{Q}}\bar y_{i,j_i^\star}.
\end{equation}
Accuracy, precision, recall, and F1 instead evaluate every
query--candidate suitability pair. These metrics need not move together:
misclassifying one of several correct candidates may leave EM unchanged, while
one highly ranked false positive can reduce EM despite otherwise accurate
pairwise predictions.

Table~\ref{tab:full-ablation} restores the precision and recall columns for the
component ablation summarized in the main paper.

\begin{table}[t]
\centering
\scriptsize
\setlength{\tabcolsep}{2.8pt}
\begin{tabular}{lccccc}
\toprule
Variant & Acc. & Prec. & Rec. & F1 & EM\\
\midrule
Freeze image
& .6542 & .6270 & .6644 & .6452 & .5351\\
Freeze text
& .5875 & .3558 & .6630 & .4631 & .3480\\
Concatenation
& .6577 & .6091 & .5979 & .6036 & .5272\\
Model name only
& .6254 & .5994 & .6351 & .6168 & .5115\\
ARMS
& \textbf{.6751} & \textbf{.6403} & \textbf{.6882}
& \textbf{.6634} & \textbf{.5721}\\
\bottomrule
\end{tabular}
\caption{Complete component ablation metrics on the four-candidate ID split.
``Model name only'' replaces the structured descriptor by the model name.}
\label{tab:full-ablation}
\end{table}

\paragraph{Ablation configurations.}
``Freeze image'' keeps CLIP fixed while updating the remaining router.
``Freeze text'' fixes DistilBERT, which affects both question and descriptor
encoding. ``Concatenation'' replaces the gated cross-attention experts with a
direct image--text feature concatenation. ``Model name only'' retains only the
candidate name as $M_j$. All variants keep the evaluation rule unchanged, so
the EM differences isolate representation and fusion choices rather than a
different selection policy.

\section{Release and Reproducibility Commitments}

Upon publication, we will release routing records, splits, normalized labels,
prompts, candidate descriptors, preprocessing, training and evaluation code,
final configurations, and an environment specification. Non-redistributable
third-party content will be represented by identifiers and reconstruction
scripts. Only metadata verified against original logs will be included;
unrecovered exact seed values, software versions, and statistical tests are not
claimed.

\paragraph{Computing resources and repeated runs.}
Router training used four NVIDIA RTX 4090 GPUs and a 16-core Intel Xeon Gold
6430 CPU; open-source VLM outputs used four NVIDIA A100 GPUs. Each configuration
was run three times with distinct random seeds, and reported metrics are the
arithmetic means; all other settings were fixed across repetitions.

\end{document}
